% 29 — Script §6.8 and first conclusion paragraph.
\sectiontag{5}\source{Use-case summary by presenters.}
\begin{frame}[t]{One search pipeline supports three kinds of tests}
\bodyfont
{\compactfont\renewcommand{\arraystretch}{1.8}
\begin{tabularx}{\linewidth}{@{}p{170pt} L L L@{}}
\toprule
Test & Reference & Target & What it reveals\\\midrule
UC1: weakness & Capacity or another protocol & Protocol under test & Underutilization and implementation bugs\\
UC2: regression & Older version & Newer version & Harmful effects of a change\\
UC3: design goal & One path & Two paths & Violations of a multipath objective\\\bottomrule
\end{tabularx}}
\vspace{24pt}
\takeaway{Choosing the reference turns a design question into an optimizable score.}
\end{frame}

% 30 — Script conclusion.
\source{}
\begin{frame}[t]{Turn hidden weaknesses into reproducible test cases}
\bodyfont
\key{AdvNet automates the search for performance problems in real protocol code.}
\vspace{28pt}
\begin{columns}[T,onlytextwidth]
\column{.47\textwidth}\lead{Current scope}
\begin{points}
\item Limited network parameters and simple traffic workloads
\item Manual root-cause investigation
\item No estimate of deployment frequency
\end{points}
\column{.47\textwidth}\lead{Next steps}
\begin{points}
\item Richer environments and workloads
\item Automatic trace simplification
\item Other domains: adaptive bitrate streaming prototype
\end{points}
\end{columns}
\end{frame}

% B1 — Script backup notes: related approaches.
\sectiontag{6}\source{}
\begin{frame}[t]{Related approaches: formal verification and GANs}
\bodyfont
\begin{columns}[T,onlytextwidth]
\column{.47\textwidth}\lead{Formal verification}
Checks specified properties using methods such as model checking and symbolic analysis.
\vspace{18pt}\par
Requires a precise model or specification; complex kernel code and changing networks create scale challenges.
\column{.47\textwidth}\lead{Generative adversarial networks}
Learn to generate data resembling a training dataset.
\vspace{18pt}\par
For AdvNet, a dataset of challenging environments is not already available.
\end{columns}
\vspace{28pt}
\takeaway{AdvNet generates candidates using measured performance gaps as feedback.}
\end{frame}

% B2 — Script backup notes: RoundRobin and OCBA.
\source{PLS comparison arranged by presenters, based on §3.3.}
\begin{frame}[t]{Alternative post-learning selection methods}
\bodyfont
{\compactfont\renewcommand{\arraystretch}{1.6}
\begin{tabularx}{\linewidth}{@{}p{170pt} L L@{}}
\toprule
Method & Allocation rule & Selection idea\\\midrule
RoundRobin & Equal budget per candidate & Return the highest observed mean\\
OCBA & More runs for uncertain candidates and close contenders & Increase the probability of selecting the best true mean\\
MRE & Remove roughly the lower half each round; stop at $\leq 5$ & Concentrate measurements on promising candidates\\\bottomrule
\end{tabularx}}
\vspace{22pt}
\textbf{OCBA:} Optimal Computing Budget Allocation
\end{frame}

% B3 — Script backup notes: GA operator equations.
\source{Operators from §3.2.2; notation by presenters.}
\begin{frame}[t]{GA operators in equations}
\bodyfont
Parents $a,b$; gene index $i$; cut points $k_1<k_2$; per-gene bounds $[l_i,u_i]$.
\vspace{12pt}
\lead{Two-point crossover}
{\fontsize{28}{35}\selectfont
\[c_i=\begin{cases}b_i & k_1<i\leq k_2,\\a_i & \text{otherwise.}\end{cases}\]}
\lead{Uniform integer mutation}
{\fontsize{28}{35}\selectfont
\[m_i=\begin{cases}U_i & \text{with probability }p,\\c_i & \text{with probability }1-p,\end{cases}
\quad U_i\sim\operatorname{UniformInteger}(l_i,u_i).\]}
\end{frame}

% B4 — Script backup notes: Fig. 11.
\source{Fig. 11 (original).}
\begin{frame}[t]{Selection methods in a statistical simulation}
\bodyfont
Vertical axis: true mean of the selected alternative.
\par\vspace{10pt}
\begin{columns}[T,onlytextwidth]
\column{555pt}\includegraphics[width=555pt]{fig11_selection.pdf}
\column{300pt}
\lead{50 Gaussian alternatives}
True means drawn uniformly from 0 to 100.
\vspace{16pt}\par
Results averaged over 2,000 trials.
\vspace{16pt}\par
\key{MRE performs best.}
\vspace{16pt}\par
{\compactfont Below 100 samples, OCBA has too few observations to estimate variance reliably.}
\end{columns}
\end{frame}

